\documentclass[10pt,a4paper]{amsart}
\usepackage[margin=19mm]{geometry}
\usepackage{amsmath,amssymb,mathtools}
\usepackage[scr=rsfs,cal=euler]{mathalfa}
\usepackage{xfrac}
\usepackage[unicode,hidelinks,bookmarks=false]{hyperref}
\newcommand{\ie}{i.e.\ }
\newcommand{\cf}{cf.\ }
\newcommand{\resp}{respectively\ }
\newcommand{\Subsection}{\subsection}
\providecommand{\nobreakdash}{\nobreak\hbox{-}}
\setlength{\emergencystretch}{2em}
\multlinegap0pt
\usepackage{amssymb}
\usepackage{tikz}
\usepackage{esint}
\usepackage{mathtools}
\usepackage{cancel}
\usepackage{caption}
\usepackage{float}
\usepackage{xcolor}
\usepackage[shortlabels]{enumitem}
\usepackage[all]{xy}
\newtheorem{lemma}{Lemma}
\newtheorem{corollary}{Corollary}
\usepackage{cleveref}
\crefname{lemma}{Lemma}{Lemmas}
\crefname{corollary}{Corollary}{Corollarys}
\newcommand{\IC}{\mathbb{C}}
\newcommand{\cl}{\mathrm{cl}}
\title{\textbf{\large{Arithmetic Deformation of Line Bundles}}}
\author{David Urbanik, Ziquan Yang}
\date{Source: arXiv:2310.08652v2 (CC BY 4.0)}
\begin{document}
\maketitle

\noindent\emph{Source and licence.} \textbf{\large{Arithmetic Deformation of Line Bundles}}. Version arXiv:2310.08652v2 (CC BY 4.0), distributed by arXiv under the Creative Commons Attribution 4.0 International licence, http://creativecommons.org/licenses/by/4.0/, as recorded in the arXiv licence metadata for this exact version and shown on the abstract page https://arxiv.org/abs/2310.08652v2. Full licence notice: \texttt{licenses/arxiv-2310.08652-CC-BY-4.0.txt}.\par\smallskip

\noindent\textit{Editorial scope.} Counted unit: the article's corollary intersection (source0.tex lines 1115--1121). The article's complete original proof of it is delivered with it (source0.tex lines 1122--1130, one \texttt{proof} environment of 9 lines). No mathematics is written, repaired or re-derived: the statement and the proof are the article's own lines, in the article's own notation, and no retained line is substituted or re-flowed.  Retained uncounted as the source-local context the proof needs: lines 1087--1093.  Nothing was omitted from any retained span, so the package carries no omission record.  No retained line contains a citation, so the wrapper carries no bibliography and no bibliography entry.  No retained line contains a figure environment, an image-inclusion command or drawing code, so no image is delivered and the package declares no asset dependency.  Editorial formatting only, in addition: an AMS \texttt{amsart} preamble that restates the article's own declarations and macros used by the retained lines, namely the environments \texttt{lemma}, \texttt{corollary} on separate counters, together with the standard packages those lines need.  The retained environments print in this document as Lemma 1, Corollary 1 in order of appearance, whereas the article prints the same items with the numbering of its own full document.  The complete native source of the article as distributed in the arXiv e-print for this exact version is bundled unchanged as \texttt{sources/148414/source0.tex}.
\begin{lemma}
\label{lem: infinite intersection}
Let $S$ be a $\IC$-variety and let $J_0 \supseteq J_1 \supseteq \cdots$ be an infinite sequence of descending constructible subsets of $S(\IC)$. Use $\cl(-)$ to denote the closure in the Zariski topology of $S(\IC)$.\footnote{We sometimes use implicitly the fact that for a constructible subset of $\IC$-points on a $\IC$-variety, analytic closure and Zariski closure coincide.} Then 
\begin{equation}
    \cl(\cap_{r \ge 0} J_r) = \cap_{r \ge 0} \cl(J_r)
\end{equation}
\end{lemma}

\begin{corollary}
\label{cor: project infinite intersection}
    Let $f : S \to T$ be a morphism between $\IC$-varieties. Let $J_0 \supseteq J_1 \supseteq \cdots$ be a descending sequence of constructible subsets of $S(\IC)$. Then 
    \begin{equation}
        \cl(f(\cap_{r \ge 0} J_r)) = \cap_{r \ge 0} \cl(f(J_r))
    \end{equation}
\end{corollary}

\begin{proof}
    Let us start by recalling a tautology: By continuity, for every subset $A \subseteq S(\IC)$, we have $f(\cl(A)) \subseteq \cl(f(A))$, or equivalently $\cl(f(\cl(A))) = \cl(f(A))$. Now we note that for $r_0 \gg 0$, 
    \begin{align*}
        \cl(f (\cap_{r \ge 0} J_r) ) &\supseteq f(\cl(\cap_{r \ge 0} J_r)) \tag{\textrm{by continuity of $f$}} \\
        &= f(\cap_{r \ge 0} \cl(J_r)) \tag{\textrm{by \cref{lem: infinite intersection}}} \\
        &=  f(\cl(J_{r_0})) \tag{\textrm{by Noetherianity of $S$}} 
    \end{align*}
    As $\cl(f (\cap_{r \ge 0} J_r))$ is closed, we moreover have $\cl(f (\cap_{r \ge 0} J_r)) \supseteq \cl(f(\cl(J_{r_0}))) = \cl(f(J_{r_0}))$. On the other hand, by Noetherianity of $T(\IC)$, we have $\cap_{r \ge 0} \cl(f(J_r)) =  \cl(f(J_{r_0}))$ for $r_0 \gg 0$. Hence we obtain $\cl(f (\cap_{r \ge 0} J_r)) \supseteq \cap_{r \ge 0} \cl(f(J_r))$. The reverse inclusion is tautological: for every $r'$, $\cl(f(\cap_{r \ge 0} J_r)) \subseteq \cl(f(J_{r'}))$ and hence $\cl(f(\cap_{r \ge 0} J_r)) \subseteq \cap_{r'} \cl(f(J_{r'}))$. 
\end{proof}

\end{document}
